% ============================================================================
% Korean Resume Template (한국형 이력서 템플릿)
% Based on: changh95/latex_resume_template_kor (GitHub 398⭐)
% Engine: XeLaTeX + kotex
% Style: Single-column, dense layout with section separators
% ============================================================================

\documentclass[a4paper, 11pt]{article}

% ============================================================================
% PACKAGES
% ============================================================================

% --- Korean Support ---
\usepackage{kotex}              % Hangul typography (auto-loads xetexko under XeLaTeX)
\usepackage{fontspec}           % XeLaTeX font selection

% --- Page Layout ---
\usepackage[
  a4paper,
  top=15mm, bottom=15mm,
  left=17mm, right=17mm
]{geometry}

% --- Typography & Sections ---
\usepackage{titlesec}           % Section header styling
\usepackage{enumitem}           % Compact list control

% --- Tables ---
\usepackage{tabularx}           % Auto-width columns for header/info
\usepackage{array}              % Column type extensions

% --- Colors & Links ---
\usepackage{xcolor}             % Color definitions
\usepackage[hidelinks]{hyperref}  % Clickable links (no colored boxes)

% --- Miscellaneous ---
\usepackage{fancyhdr}           % Header/footer control
\usepackage{calc}               % Length arithmetic

% ============================================================================
% FONTS
% ============================================================================

\setmainfont{Malgun Gothic}[
  Script = Hangul,
  Language = Korean
]

\setmainhangulfont{Malgun Gothic}
\setsanshangulfont{Malgun Gothic}

% ============================================================================
% PAGE SETUP
% ============================================================================

\geometry{a4paper, top=15mm, bottom=15mm, left=17mm, right=17mm}

% Disable default page style
\pagestyle{empty}

% Compact spacing throughout
\setlength{\parindent}{0pt}
\setlength{\parskip}{0pt}
\setlength{\baselineskip}{1.1\baselineskip}

% ============================================================================
% COLORS
% ============================================================================

\definecolor{navy}{RGB}{31, 73, 125}        % Dark blue for accents
\definecolor{darkgray}{RGB}{60, 60, 60}     % Dark gray for section headers

% ============================================================================
% SECTION FORMATTING (titlesec)
% ============================================================================

% Section: Bold title + full-width black line below
\titleformat{\section}
  {\color{darkgray}\large\bfseries}  % Format
  {}                                 % Label (none)
  {0pt}                              % Separation
  {}                                 % Before
  [\vspace{-2pt}\color{darkgray}\hrule\vspace{4pt}]  % After: thin line + spacing

\titlespacing{\section}{0pt}{6pt plus 2pt minus 2pt}{4pt plus 1pt minus 1pt}

% ============================================================================
% LISTS
% ============================================================================

% Compact bullet lists
\setlist[itemize]{
  noitemsep,        % No extra space between items
  topsep=0pt,       % No space above list
  partopsep=0pt,    % No space if list is first in paragraph
  leftmargin=12pt   % Minimal indentation
}

% ============================================================================
% CUSTOM COMMANDS
% ============================================================================

% Header: Name + Contact Info (side-by-side in a table)
\newcommand{\resumeHeader}[5]{%
  \begin{center}
    {\LARGE\bfseries #1}
  \end{center}
  \vspace{2pt}
  \noindent
  \begin{tabularx}{\textwidth}{X r}
    & 전화: #2 \quad 이메일: \href{mailto:#3}{#3} \\
    & 주소: #4 \quad 블로그: \href{#5}{#5}
  \end{tabularx}
  \vspace{6pt}
}

% Entry with Subheading (title + metadata on two lines)
% Args: {org/school}{period}{position/major/role}{location/city}
\newcommand{\resumeSubheading}[4]{%
  \vspace{2pt}
  \begin{tabularx}{\textwidth}{l r}
    \textbf{#1} & \small #2 \\
    #3 & \small #4 \\
  \end{tabularx}
  \vspace{2pt}
}

% Compact bullet item
\newcommand{\resumeItem}[1]{%
  \item \small #1
}

% Skill category + list (single line)
\newcommand{\resumeSkill}[2]{%
  \textbf{#1:} #2 \\
}

% ============================================================================
% DOCUMENT START
% ============================================================================

\begin{document}

% ============================================================================
% HEADER: NAME + CONTACT
% ============================================================================

\resumeHeader
  {조 동 휘}  % Name (with spacing)
  {010-8699-1510}  % Phone
  {donghwi0711@icloud.com}  % Email
  {서울특별시 관악구 신림로 30길 17-19 (올리하우스) 306호}  % Address
  {https://github.com/nobrain711}  % Blog/Portfolio URL

% ============================================================================
% SECTION: 학력사항 (Education)
% ============================================================================

\section{학력사항}

\resumeSubheading
  {영남이공대학교}
  {2021.03 -- 2024.02}
  {컴퓨터공학과 (학점: 3.75/4.5)}
  {대구광역시}

\resumeSubheading
  {춘천기계공업고등학교}
  {2016.03 -- 2019.02}
  {특성화고등학교 자동차과(군특성화과)}
  {강원도 춘천시}

% ============================================================================
% SECTION: 경력사항 (Work Experience)
% ============================================================================

\section{경력사항}

\resumeSubheading
  {(주)토마토 (일본 개발 인력 파견회사)}
  {2024.04 -- 2025.01}
  {소프트웨어 엔지니어}
  {일본 도쿄도 시부야구}

\begin{itemize}
  \resumeItem{금융계 시스템 개발 담당 — AWS EC2 기반 클라우드 서버 및 온프레미스 물리 서버 환경 구축}
  \resumeItem{Windows Server 19/21, Microsoft SQL Server 19/21 설치 및 환경 설정 (클라우드 \& 온프레미스 이중 구성)}
  \resumeItem{Visual Studio 설치 작업 지시서(SOP) 작성 — 과거 버전 대비 명확성 개선으로 설치 시간 14\% 단축 달성}
  \resumeItem{금융 시스템 운영 및 유지보수 경험 (6개월, 파견사원)}
\end{itemize}

% ============================================================================
% SECTION: 기술스택 (Technical Skills)
% ============================================================================

\section{기술스택}

\noindent
\resumeSkill{프로그래밍 언어}{Python, JavaScript/TypeScript, SQL, Bash}

\noindent
\resumeSkill{웹 프레임워크}{FastAPI, Django, React, Express.js}

\noindent
\resumeSkill{데이터베이스}{PostgreSQL, MySQL, MongoDB, Redis}

\noindent
\resumeSkill{클라우드 \& DevOps}{AWS (EC2, S3, RDS), Docker, GitHub Actions, Kubernetes}

\noindent
\resumeSkill{개발 도구}{Git, VS Code, Postman, Jupyter Notebook}

\noindent
\resumeSkill{기타 라이브러리}{SQLAlchemy, Pydantic, NumPy, Pandas, Matplotlib}

% ============================================================================
% SECTION: 자격증 / 어학 (Certifications & Language)
% ============================================================================

\section{자격증 \& 어학}

\begin{tabularx}{\textwidth}{l l l}
  \textbf{자격증} & & \\
  정보처리기사 & 한국산업인력공단 & 2021.11.26 \\
  SQLD & 한국데이터산업진흥원 & 2022.04.15 \\
  & & \\
  \textbf{어학} & & \\
  TOEIC & 920점 & 2023.05.14 \\
  OPIc & IM2 & 2023.09.10 \\
\end{tabularx}

% ============================================================================
% SECTION: 수상 / 활동 (Awards & Activities)
% ============================================================================

\section{수상 \& 활동}

\resumeSubheading
  {교내 프로그래밍 경진대회}
  {2019.11}
  {최우수상 (1위)}
  {한국대학교}

\resumeSubheading
  {오픈소스 프로젝트 컨트리뷰션}
  {2022.06 -- 현재}
  {GitHub 저장소 기여 (40+ commits)}
  {온라인}

\resumeSubheading
  {지역 IT 교육 봉사}
  {2021.03 -- 2021.08}
  {청소년 대상 코딩 교육 지도 (월 4회)}
  {서울시 강남구}

\end{document}
